\section*{第 \ref{ch:g6:earth-a-planet} 章——地球，太阳系中的一颗行星}

\begin{solution}{exo:g6:earth-a-planet:1}
\underline{水星}、\underline{金星}、\underline{地球}、
\underline{火星}、木星、土星、天王星、海王星。\omterm{def:g6:earth-a-planet:planet-types}{岩质行星}是
前四颗。
\end{solution}

\begin{solution}{exo:g6:earth-a-planet:2}
“\omterm{def:g6:earth-a-planet:planet-types}{岩质行星}体积小、密度大，表面为固体。\omterm{def:g6:earth-a-planet:planet-types}{巨行星}
体积非常大，密度小，主要由气体和液体组成。”
\end{solution}

\begin{solution}{exo:g6:earth-a-planet:3}
金星（12\,104~千米，地球为 12\,756~千米）。最大的是
木星（142\,984~千米）；最小的是水星（4\,879~千米）。
\end{solution}

\begin{solution}{exo:g6:earth-a-planet:4}
（a）错误：金星上没有液态水，温度太高。（b）正确
（平均 $-65$~$^\circ$C，地球为 15~$^\circ$C）。（c）错误：
离太阳最近的四颗行星是\omterm{def:g6:earth-a-planet:planet-types}{岩质行星}；\omterm{def:g6:earth-a-planet:planet-types}{巨行星}是
离太阳最远的四颗。（d）正确：约为 0.7~克每立方厘米。
\end{solution}

\begin{solution}{exo:g6:earth-a-planet:5}
$4879 \div 12756 \approx 0.382$，即 38~\%：水星的直径略多于
地球直径的三分之一。
\end{solution}

\begin{solution}{exo:g6:earth-a-planet:6}
$120536 \div 12756 \approx 9.4$：约为 9 个地球。土星是\omterm{def:g6:earth-a-planet:planet-types}{巨行星}。
\end{solution}

\begin{solution}{exo:g6:earth-a-planet:7}
行星~A 是\omterm{def:g6:earth-a-planet:planet-types}{巨行星}：它体积很大，其密度为
1.1~克每立方厘米，很低，与太阳系\omterm{def:g6:earth-a-planet:planet-types}{巨行星}的密度相近，
（介于 0.7 和 1.7 之间）。行星~B 是\omterm{def:g6:earth-a-planet:planet-types}{岩质行星}：它体积小，
密度为 5.0，较高，与\omterm{def:g6:earth-a-planet:planet-types}{岩质行星}的密度相近（介于 3.9 和
5.5 之间）。
\end{solution}

\begin{solution}{exo:g6:earth-a-planet:8}
（a）$464 - 15 = 449$~度更热。（b）$15 - (-65) = 80$~度更冷。
\end{solution}

\begin{solution}{exo:g6:earth-a-planet:9}
金星拥有非常厚的二氧化碳大气，像一条厚毯子一样困住太阳的
热量，使其无法散失。水星离太阳更近，
几乎没有空气，因此无法保存热量。
\end{solution}

\begin{solution}{exo:g6:earth-a-planet:10}
$228 - 150 = 78$~百万千米更远。$78 \div 150 = 0.52$：火星
与太阳的距离比地球的大 52~\%。
\end{solution}

\begin{solution}{exo:g6:earth-a-planet:11}
数十亿年前，火星上有\omterm{def:g1:landscapes-around-us:water}{河流}、\omterm{def:g1:landscapes-around-us:water}{湖泊}和\omterm{def:g5:water-shapes-the-land:flood}{洪水}：那时它更温暖、
更湿润。如今，其水以冰的形式存在于极冠中，并埋在
靠近两极的地下。如今火星表面无法长期存在液态水，
因为火星非常寒冷（平均 $-65$~$^\circ$C），且其空气
极为稀薄，水很快就会凝固或变成蒸气。
\end{solution}

\begin{solution}{exo:g6:earth-a-planet:12}
（a）$8.9 \div 9.8 \approx 0.908$，即 91~\%。（b）相似之处：几乎
大小相同（直径为地球的 95~\%），引力也几乎相同，
它也是\omterm{def:g6:earth-a-planet:planet-types}{岩质行星}。不同之处：它的温度高得多
（464~$^\circ$C，地球为 15~$^\circ$C），大气厚得多，由二氧化碳构成，
并有酸性\omterm{def:g4:clouds-rain-snow:cloud}{云层}，而且没有液态水。
\end{solution}

\begin{solution}{pb:g6:earth-a-planet:1}
\textbf{1.} 金星（108~百万千米）、地球（150）、火星（228）。

\textbf{2.} $150 - 108 = 42$：金星比地球近 42~百万千米。
$228 - 150 = 78$：火星比地球远 78~百万千米。

\textbf{3.} $6792 \div 12756 \approx 0.532$，即 53~\%：火星的直径约为
地球的一半。

\textbf{4.} 三者都是\omterm{def:g6:earth-a-planet:planet-types}{岩质行星}：它们的密度（分别为 5.2、5.5 和 3.9
克每立方厘米）均较高，直径
与\omterm{def:g6:earth-a-planet:planet-types}{巨行星}相比较小。

\textbf{5.} 金星的大气最厚（92~巴），火星的最薄
（0.01~巴）。

\textbf{6.} $1 \div 0.01 = 100$：火星的气压
比地球的小一百倍。

\textbf{7.} 其厚厚的二氧化碳大气困住太阳的热量，使其无法散失，
像一床厚毯子。

\textbf{8.} $92 \div 1 = 92$：在金星地面，气压
是地球地面的 \textbf{92 倍}。探测器必须有非常坚固的外壳，
才不会被压坏（还必须能抵抗高温
和酸性\omterm{def:g4:clouds-rain-snow:cloud}{云层}）。

\textbf{9.} 地球的平均温度为 15~$^\circ$C，介于
0~$^\circ$C 与 100~$^\circ$C 之间：水既不结冰也不沸腾，因此
保持液态。

\textbf{10.} 在金星上，温度为 464~$^\circ$C，远高于 100~$^\circ$C，水
会立即沸腾并变成蒸气。在火星上，温度为 $-65$~$^\circ$C，
远低于 0~$^\circ$C，水会凝固成冰（并因空气极为稀薄而缓慢
变成蒸气）。

\textbf{11.} 很久以前，火星比今天更温暖、更湿润：其表面
可有液态水流动。其\omterm{def:g5:weather-and-climate:climate}{气候}已发生很大变化：如今它是一片
寒冷干燥的\omterm{def:g5:weather-and-climate:desert}{沙漠}。

\textbf{12.} 例如：“亲爱的探测器们，我们地球上有液态水，
存在于\omterm{def:g2:water-on-earth:ocean}{海洋}、\omterm{def:g1:landscapes-around-us:water}{河流}和\omterm{def:g1:landscapes-around-us:water}{湖泊}中。我们的气温适中，平均为 15~$^\circ$C，
介于水的凝固与沸腾之间。我们与太阳的距离
恰到好处，既不太近，也不太远。我们的大气
使我们保持温暖，为生物提供所需的气体，并让
水保持液态。这就是这里有生命存在的原因。”
\end{solution}
